% Fichier généré automatiquement par tools/generate_corrections.py.
% Source : GEO/GEO-01/14_Sympact/14_Sympact.tex
% Statut : complété (5/5 question(s)).
\begin{corrigebox}[Corrigé complété]
\CorrigeQuestion{1}
\begin{center}
\includegraphics[width=.4\linewidth]{14_02_c}
\end{center}
\CorrigeQuestion{2}
\begin{center}
\includegraphics[width=.5\linewidth]{14_01_c}
\end{center}
\CorrigeQuestion{3}
On reconstruit la configuration en appliquant les
                rotations dans l'ordre de la chaîne cinématique puis les translations le
                long de leurs axes. Les longueurs des barres restent constantes. Les
                coordonnées finales se calculent avec
                $\vec i'=\cos q\,\vec i+\sin q\,\vec j$; le tracé doit conserver les
                liaisons et le contact sans pénétration.
\CorrigeQuestion{4}
Dans cas, dans le pire des cas, $A$, $B$ et $C$ sont alignés (avec $B$ au-dessus de $C$). Il faut donc $AB = AC+CB = \SI{160}{mm}$.
\CorrigeQuestion{5}
Comme je suis paresseux, j'ai réalisé la construction avec geogebra. On mesure \SI{160,8}{\degres}.
\begin{marginfigure}
\rotatebox{90}{\includegraphics[width=.35\linewidth]{14_03_c}}
\end{marginfigure}

\begin{enumerate}
\item .
\item .
\item .
\item \SI{160}{mm}.
\item \SI{160,8}{\degres}.
\end{enumerate}
\end{corrigebox}
